% Generated by tools/build_prior_work.py from the claims ledger; do not edit.
\bibitem{pwaharonov1959significance}
Aharonov, Y., and Bohm, D. (1959). ``Significance of electromagnetic potentials in the quantum theory.'' \emph{Physical Review} \textbf{115}, 485--491. \url{https://doi.org/10.1103/PhysRev.115.485}

\bibitem{pwbaez2010algebra}
Baez, J., and Huerta, J. (2010). ``The algebra of grand unified theories.'' \emph{Bulletin of the American Mathematical Society} \textbf{47}, 483--552. \url{https://doi.org/10.1090/S0273-0979-10-01294-2}

\bibitem{pwbarndorffnielsen2014information}
Barndorff-Nielsen, O. (2014). \emph{Information and Exponential Families: In Statistical Theory}. Wiley (Wiley Series in Probability and Statistics). \url{https://doi.org/10.1002/9781118857281}

\bibitem{pwbarnich2000local}
Barnich, G., Brandt, F., and Henneaux, M. (2000). ``Local BRST cohomology in gauge theories.'' \emph{Physics Reports} \textbf{338}, 439--569. \url{https://doi.org/10.1016/S0370-1573(00)00049-1}

\bibitem{pwbecchi1976renormalization}
Becchi, C., Rouet, A., and Stora, R. (1976). ``Renormalization of gauge theories.'' \emph{Annals of Physics} \textbf{98}, 287--321. \url{https://doi.org/10.1016/0003-4916(76)90156-1}

\bibitem{pwbouchiat1972anomalyfree}
Bouchiat, C., Iliopoulos, J., and Meyer, P. (1972). ``An anomaly-free version of Weinberg's model.'' \emph{Physics Letters B} \textbf{38}, 519--523. \url{https://doi.org/10.1016/0370-2693(72)90532-1}

\bibitem{pwboyd2004convex}
Boyd, S., and Vandenberghe, L. (2004). \emph{Convex Optimization}. Cambridge University Press. \url{https://doi.org/10.1017/CBO9780511804441}

\bibitem{pwbratteli1997operator}
Bratteli, O., and Robinson, D. W. (1997). \emph{Operator Algebras and Quantum Statistical Mechanics 2}. Springer. \url{https://doi.org/10.1007/978-3-662-03444-6}

\bibitem{pwbrocker1985representations}
Br{\"{o}}cker, T., and tom Dieck, T. (1985). \emph{Representations of Compact Lie Groups}. Springer. \url{https://doi.org/10.1007/978-3-662-12918-0}

\bibitem{pwbrouwer1989distance}
Brouwer, A. E., Cohen, A. M., and Neumaier, A. (1989). \emph{Distance-Regular Graphs}. Springer, Berlin. \url{https://doi.org/10.1007/978-3-642-74341-2}

\bibitem{pwcohn2007universally}
Cohn, H., and Kumar, A. (2007). ``Universally optimal distribution of points on spheres.'' \emph{Journal of the American Mathematical Society} \textbf{20}, 99--148. \url{https://doi.org/10.1090/S0894-0347-06-00546-7}

\bibitem{pwcoxeter1973regular}
Coxeter, H. S. M. (1973). \emph{Regular Polytopes}. Dover. ISBN 0-486-61480-8.

\bibitem{pwcoxeter1980generators}
Coxeter, H. S. M., and Moser, W. O. J. (1980). \emph{Generators and Relations for Discrete Groups}. Springer Berlin Heidelberg. \url{https://doi.org/10.1007/978-3-662-21943-0}

\bibitem{pwdavies1970operational}
Davies, E. B., and Lewis, J. T. (1970). ``An operational approach to quantum probability.'' \emph{Communications in Mathematical Physics} \textbf{17}, 239--260. \url{https://doi.org/10.1007/BF01647093}

\bibitem{pwdeligne1982tannakian}
Deligne, P., and Milne, J. S. (1982). ``Tannakian Categories.'' In \emph{Hodge Cycles, Motives, and Shimura Varieties}, 101--228. Springer Berlin Heidelberg. \url{https://doi.org/10.1007/978-3-540-38955-2_4}

\bibitem{pwdeligne2002categories}
Deligne, P. (2002). ``Cat{\'{e}}gories tensorielles.'' \emph{Moscow Mathematical Journal} \textbf{2}, 227--248. \url{https://doi.org/10.17323/1609-4514-2002-2-2-227-248}

\bibitem{pwdenner1993techniques}
Denner, A. (1993). ``Techniques for the calculation of electroweak radiative corrections at the one-loop level and results for W-physics at LEP 200.'' \emph{Fortschritte der Physik} \textbf{41}, 307--420. \url{https://doi.org/10.1002/prop.2190410402}

\bibitem{pwdiaconis1988group}
Diaconis, P. (1988). \emph{Group Representations in Probability and Statistics}. Institute of Mathematical Statistics (Lecture Notes--Monograph Series). \url{https://doi.org/10.1214/lnms/1215467407}

\bibitem{pwdoplicher1971local}
Doplicher, S., Haag, R., and Roberts, J. E. (1971). ``Local observables and particle statistics I.'' \emph{Communications in Mathematical Physics} \textbf{23}, 199--230. \url{https://doi.org/10.1007/BF01877742}

\bibitem{pwdoplicher1989new}
Doplicher, S., and Roberts, J. E. (1989). ``A new duality theory for compact groups.'' \emph{Inventiones Mathematicae} \textbf{98}, 157--218. \url{https://doi.org/10.1007/BF01388849}

\bibitem{pwdoplicher1990why}
Doplicher, S., and Roberts, J. E. (1990). ``Why there is a field algebra with a compact gauge group describing the superselection structure in particle physics.'' \emph{Communications in Mathematical Physics} \textbf{131}, 51--107. \url{https://doi.org/10.1007/BF02097680}

\bibitem{pweden1966analytic}
Eden, R. J., Landshoff, P. V., Olive, D. I., and Polkinghorne, J. C. (1966). \emph{The Analytic S-Matrix}. Cambridge University Press. ISBN 0-521-04869-9.

\bibitem{pwfleischer1981radiative}
Fleischer, J., and Jegerlehner, F. (1981). ``Radiative corrections to Higgs-boson decays in the Weinberg-Salam model.'' \emph{Physical Review D} \textbf{23}, 2001--2026. \url{https://doi.org/10.1103/PhysRevD.23.2001}

\bibitem{pwfock1932konfigurationsraum}
Fock, V. (1932). ``Konfigurationsraum und zweite Quantelung.'' \emph{Zeitschrift f{\"{u}}r Physik} \textbf{75}, 622--647. \url{https://doi.org/10.1007/BF01344458}

\bibitem{pwfulton1991representation}
Fulton, W., and Harris, J. (1991). \emph{Representation Theory: A First Course}. Springer, New York (Graduate Texts in Mathematics). \url{https://doi.org/10.1007/978-1-4612-0979-9}

\bibitem{pwgambino2000nielsen}
Gambino, P., and Grassi, P. A. (2000). ``Nielsen identities of the SM and the definition of mass.'' \emph{Physical Review D} \textbf{62}, 076002. \url{https://doi.org/10.1103/PhysRevD.62.076002}

\bibitem{pwgeorgi1974unity}
Georgi, H., and Glashow, S. L. (1974). ``Unity of All Elementary-Particle Forces.'' \emph{Physical Review Letters} \textbf{32}, 438--441. \url{https://doi.org/10.1103/PhysRevLett.32.438}

\bibitem{pwginsparg1982remnant}
Ginsparg, P. H., and Wilson, K. G. (1982). ``A remnant of chiral symmetry on the lattice.'' \emph{Physical Review D} \textbf{25}, 2649--2657. \url{https://doi.org/10.1103/PhysRevD.25.2649}

\bibitem{pwgreenberg1991particles}
Greenberg, O. W. (1991). ``Particles with small violations of Fermi or Bose statistics.'' \emph{Physical Review D} \textbf{43}, 4111--4120. \url{https://doi.org/10.1103/PhysRevD.43.4111}

\bibitem{pwgross1972effect}
Gross, D. J., and Jackiw, R. (1972). ``Effect of Anomalies on Quasi-Renormalizable Theories.'' \emph{Physical Review D} \textbf{6}, 477--493. \url{https://doi.org/10.1103/PhysRevD.6.477}

\bibitem{pwgrzadkowski2010dimensionsix}
Grzadkowski, B., Iskrzy{\'{n}}ski, M., Misiak, M., and Rosiek, J. (2010). ``Dimension-six terms in the Standard Model Lagrangian.'' \emph{Journal of High Energy Physics} \textbf{2010}, 85. \url{https://doi.org/10.1007/JHEP10(2010)085}

\bibitem{pwhalmos1958finite}
Halmos, P. R. (1958). \emph{Finite-Dimensional Vector Spaces}. Van Nostrand. \url{https://doi.org/10.1007/978-1-4612-6387-6}

\bibitem{pwhartman2002ordinary}
Hartman, P. (2002). \emph{Ordinary Differential Equations}. Society for Industrial and Applied Mathematics (Classics in Applied Mathematics). \url{https://doi.org/10.1137/1.9780898719222}

\bibitem{pwhenneaux1992quantization}
Henneaux, M., and Teitelboim, C. (1992). \emph{Quantization of Gauge Systems}. Princeton University Press. \url{https://doi.org/10.1515/9780691213866}

\bibitem{pwhong1987measurement}
Hong, C. K., Ou, Z. Y., and Mandel, L. (1987). ``Measurement of subpicosecond time intervals between two photons by interference.'' \emph{Physical Review Letters} \textbf{59}, 2044--2046. \url{https://doi.org/10.1103/PhysRevLett.59.2044}

\bibitem{pwhucks1991global}
Hucks, J. (1991). ``Global structure of the standard model, anomalies, and charge quantization.'' \emph{Physical Review D} \textbf{43}, 2709--2717. \url{https://doi.org/10.1103/PhysRevD.43.2709}

\bibitem{pwjames2001representations}
James, G., and Liebeck, M. (2001). \emph{Representations and Characters of Groups}. Cambridge University Press. \url{https://doi.org/10.1017/CBO9780511814532}

\bibitem{pwjoyal1991introduction}
Joyal, A., and Street, R. (1991). ``An introduction to Tannaka duality and quantum groups.'' In \emph{Category Theory (Como, 1990)}, 413--492. Springer Berlin Heidelberg. \url{https://doi.org/10.1007/BFb0084235}

\bibitem{pwknapp1996lie}
Knapp, A. W. (1996). \emph{Lie Groups Beyond an Introduction}. Birkh{\"{a}}user, Boston (Progress in Mathematics). \url{https://doi.org/10.1007/978-1-4757-2453-0}

\bibitem{pwkogut1975hamiltonian}
Kogut, J., and Susskind, L. (1975). ``Hamiltonian formulation of Wilson's lattice gauge theories.'' \emph{Physical Review D} \textbf{11}, 395--408. \url{https://doi.org/10.1103/PhysRevD.11.395}

\bibitem{pwkogut1979introduction}
Kogut, J. B. (1979). ``An introduction to lattice gauge theory and spin systems.'' \emph{Reviews of Modern Physics} \textbf{51}, 659--713. \url{https://doi.org/10.1103/RevModPhys.51.659}

\bibitem{pwkraus1983states}
Kraus, K. (1983). \emph{States, Effects, and Operations: Fundamental Notions of Quantum Theory}. Springer. \url{https://doi.org/10.1007/3-540-12732-1}

\bibitem{pwlangacker2017standard}
Langacker, P. (2017). \emph{The Standard Model and Beyond}. CRC Press. \url{https://doi.org/10.1201/b22175}

\bibitem{pwlemmens1973equiangular}
Lemmens, P. W. H., and Seidel, J. J. (1973). ``Equiangular lines.'' \emph{Journal of Algebra} \textbf{24}, 494--512. \url{https://doi.org/10.1016/0021-8693(73)90123-3}

\bibitem{pwloudon1998fermion}
Loudon, R. (1998). ``Fermion and boson beam-splitter statistics.'' \emph{Physical Review A} \textbf{58}, 4904--4909. \url{https://doi.org/10.1103/PhysRevA.58.4904}

\bibitem{pwluscher1999abelian}
L{\"{u}}scher, M. (1999). ``Abelian chiral gauge theories on the lattice with exact gauge invariance.'' \emph{Nuclear Physics B} \textbf{549}, 295--334. \url{https://doi.org/10.1016/S0550-3213(99)00115-7}

\bibitem{pwmilnor1997topology}
Milnor, J. W. (1997). \emph{Topology from the Differentiable Viewpoint}. Princeton University Press (Princeton Landmarks in Mathematics). ISBN 0-691-04833-9.

\bibitem{pwmontvay1994quantum}
Montvay, I., and M{\"{u}}nster, G. (1994). \emph{Quantum Fields on a Lattice}. Cambridge University Press. \url{https://doi.org/10.1017/CBO9780511470783}

\bibitem{pwmoore2009introduction}
Moore, R. E., Kearfott, R. B., and Cloud, M. J. (2009). \emph{Introduction to Interval Analysis}. SIAM. \url{https://doi.org/10.1137/1.9780898717716}

\bibitem{pwnielsen1975gauge}
Nielsen, N. K. (1975). ``On the gauge dependence of spontaneous symmetry breaking in gauge theories.'' \emph{Nuclear Physics B} \textbf{101}, 173--188. \url{https://doi.org/10.1016/0550-3213(75)90301-6}

\bibitem{pwosterwalder1973axioms}
Osterwalder, K., and Schrader, R. (1973). ``Axioms for Euclidean Green's functions.'' \emph{Communications in Mathematical Physics} \textbf{31}, 83--112. \url{https://doi.org/10.1007/BF01645738}

\bibitem{pwosterwalder1975axioms}
Osterwalder, K., and Schrader, R. (1975). ``Axioms for Euclidean Green's functions II.'' \emph{Communications in Mathematical Physics} \textbf{42}, 281--305. \url{https://doi.org/10.1007/BF01608978}

\bibitem{pwpeskin1995introduction}
Peskin, M. E., and Schroeder, D. V. (1995). \emph{An Introduction to Quantum Field Theory}. Addison-Wesley. \url{https://doi.org/10.1201/9780429503559}

\bibitem{pwpiguet1995algebraic}
Piguet, O., and Sorella, S. P. (1995). \emph{Algebraic Renormalization: Perturbative Renormalization, Symmetries and Anomalies}. Springer. \url{https://doi.org/10.1007/978-3-540-49192-7}

\bibitem{pwserre1977linear}
Serre, J. P. (1977). \emph{Linear Representations of Finite Groups}. Springer. \url{https://doi.org/10.1007/978-1-4684-9458-7}

\bibitem{pwslansky1981group}
Slansky, R. (1981). ``Group theory for unified model building.'' \emph{Physics Reports} \textbf{79}, 1--128. \url{https://doi.org/10.1016/0370-1573(81)90092-2}

\bibitem{pwstuart1991gauge}
Stuart, R. G. (1991). ``Gauge invariance, analyticity and physical observables at the Z0 resonance.'' \emph{Physics Letters B} \textbf{262}, 113--119. \url{https://doi.org/10.1016/0370-2693(91)90653-8}

\bibitem{pwtannaka1939dualitatssatz}
Tannaka, T. (1939). ``{\"{U}}ber den Dualit{\"{a}}tssatz der nichtkommutativen topologischen Gruppen.'' \emph{Tohoku Mathematical Journal, First Series} \textbf{45}, 1--12. \url{https://www.jstage.jst.go.jp/article/tmj1911/45/0/45_0_1/_article}

\bibitem{pwtong2017line}
Tong, D. (2017). ``Line operators in the Standard Model.'' \emph{Journal of High Energy Physics} \textbf{2017}, 104. \url{https://doi.org/10.1007/JHEP07(2017)104}

\bibitem{pwweinberg1967model}
Weinberg, S. (1967). ``A model of leptons.'' \emph{Physical Review Letters} \textbf{19}, 1264--1266. \url{https://doi.org/10.1103/PhysRevLett.19.1264}

\bibitem{pwweinberg1979baryon}
Weinberg, S. (1979). ``Baryon- and Lepton-Nonconserving Processes.'' \emph{Physical Review Letters} \textbf{43}, 1566--1570. \url{https://doi.org/10.1103/PhysRevLett.43.1566}

\bibitem{pwweinberg1996quantum}
Weinberg, S. (1996). \emph{The Quantum Theory of Fields, Volume II}. Cambridge University Press. \url{https://doi.org/10.1017/CBO9781139644174}

\bibitem{pwwielandt1949unbeschranktheit}
Wielandt, H. (1949). ``{\"{U}}ber die Unbeschr{\"{a}}nktheit der Operatoren der Quantenmechanik.'' \emph{Mathematische Annalen} \textbf{121}, 21. \url{https://doi.org/10.1007/BF01329611}

\bibitem{pwwilczek1979operator}
Wilczek, F., and Zee, A. (1979). ``Operator Analysis of Nucleon Decay.'' \emph{Physical Review Letters} \textbf{43}, 1571--1573. \url{https://doi.org/10.1103/PhysRevLett.43.1571}

\bibitem{pwwilson1974confinement}
Wilson, K. G. (1974). ``Confinement of quarks.'' \emph{Physical Review D} \textbf{10}, 2445--2459. \url{https://doi.org/10.1103/PhysRevD.10.2445}

\bibitem{pwyukawa1935interaction}
Yukawa, H. (1935). ``On the interaction of elementary particles. I.'' \emph{Proceedings of the Physico-Mathematical Society of Japan. 3rd Series} \textbf{17}, 48--57. \url{https://doi.org/10.11429/ppmsj1919.17.0_48}

\bibitem{pwzaslavsky1982signed}
Zaslavsky, T. (1982). ``Signed graphs.'' \emph{Discrete Applied Mathematics} \textbf{4}, 47--74. \url{https://doi.org/10.1016/0166-218X(82)90033-6}

